\section{总结与延伸}

这堂课最重要的贡献不是一份 engine 或 GPU 推荐清单，而是一条可复用的 production inference 方法论：从应用语义开始，用 workload/SLO 固化目标，在单副本上画 frontier，选择模型与 engine，把 prefill/decode 映射到硬件，再设计弹性、故障、观测和正确性，最后按数量级依次优化。

\subsection{十步部署检查表}

\begin{importantbox}{从需求到 kernel 的十步顺序}
\begin{enumerate}
  \item 用 application eval 定义 correctness 与允许损失；
  \item 记录 QPS/RPS、输入输出 token、prefix reuse、burst 与 tool loop；
  \item 声明 TTFT、TPOT/ITL、TTLT、成本与可靠性 SLO；
  \item 在单 replica 上跑 serial、burst 与 rate sweep；
  \item 判断 capability-bound 或 efficiency-bound，再选模型；
  \item 用相同 workload 比较 engine，而不是比较不同 demo；
  \item 用 arithmetic intensity、HBM、Tensor Core 与 interconnect 选硬件；
  \item 设计 allocation、health、cache、checkpoint 与 headroom；
  \item 统一记录 token IDs、queue/stage latency、replica、hardware 与 eval；
  \item 按 speculation/quantization、host、kernel 的顺序优化并 canary。
\end{enumerate}
\end{importantbox}

\subsection{最常见的因果混淆}

\begin{warningbox}{六个“看起来更快”但可能不成立的结论}
\begin{enumerate}
  \item Aggregate tokens/s 更高，不代表 tokens/s/user 或 p99 更好；
  \item GPU utilization 高，不代表 SM、HBM 或 Tensor Core 有效利用；
  \item Provider price 低，不代表你的 workload 自托管更贵或更便宜；
  \item Quantized perplexity（交叉熵的指数，可粗略理解为每个 token 的有效选择数）接近，不代表 tool use 与长上下文正确性不变；
  \item Kernel microbenchmark 加速，不代表端到端 latency 改善；
  \item Agent 能生成代码，不代表系统具备 correctness、权限和 rollback。
\end{enumerate}
\end{warningbox}

\subsection{一张端到端 mental model}

\begin{knowledgebox}{Production inference control loop}
应用与 eval 定义目标 $\rightarrow$ workload/SLO 定义负载与约束 $\rightarrow$ model/engine/hardware 给出单副本 frontier $\rightarrow$ allocator 与 scheduler 把总流量映射到 replica $\rightarrow$ logs/metrics/evals 观测质量和性能 $\rightarrow$ optimizer/agent 提出改动 $\rightarrow$ correctness gate、benchmark、canary 与 rollback 决定是否上线。任何跳过测量或验证的捷径，都会把未知风险推到用户侧。
\end{knowledgebox}

\subsection{拓展阅读}

\begin{itemize}
  \item vLLM / PagedAttention：arXiv:2309.06180；
  \item Orca iteration-level scheduling：OSDI 2022；
  \item SGLang / RadixAttention：arXiv:2312.07104；
  \item FlashAttention-2：arXiv:2307.08691；
  \item FlashInfer：arXiv:2501.01005；
  \item Lossless speculative decoding：arXiv:2211.17192；
  \item EAGLE：arXiv:2401.15077；
  \item DFlash：arXiv:2602.06036；
  \item FP8 formats：arXiv:2209.05433；
  \item VibeServe：arXiv:2605.06068；
  \item NVIDIA CUDA Graph / Nsight 官方文档；
  \item Modal LLM Almanac、GPU Glossary 与 serverless GPU 工程文章。
\end{itemize}

\begin{knowledgebox}{Lecture snapshot：2026-05-28}
本讲义冻结在课堂日期可获得的一手资料。Slides 070、072、073 未作为课堂内容重建：070 是短暂过渡笑话，072 的 CI/CL appendix 未在录像中讲到，073 是招聘页。完整录像以五秒间隔审计 990 帧、17 张 contact sheets；除 token timing simulator 外，没有 deck 外独立白板或 live coding 视觉证据。
\end{knowledgebox}

\end{document}
